\documentclass[11pt]{article}\usepackage{layout}\begin{document}
\begin{canvas}\Header{第二章 · 全球开源生态地图}{teal}\Footer{01}
\Box{15}{27}{180}{\fontsize{21}{27}\selectfont\color{ink}\sffamily \bfseries 2.3 全球开源项目影响地图}
\Box{15}{40}{180}{\fontsize{9.7}{13.5}\selectfont\color{muted}\sffamily 项目影响力以 OpenRank 衡量项目在协作网络中的位置，反映贡献关系与网络连接。它不直接衡量用户规模或商业价值；项目归属地也可能不同于贡献者所在地。}
\Box{15}{53}{180}{\includegraphics[width=180mm]{.build/assets/01-project-influence-map.pdf}}
\Box{15}{199}{180}{\fontsize{7.8}{10.8}\selectfont\color{muted}\sffamily 图 2-8  全球项目影响、参与结构与五域比较。来源：OpenDigger / OpenShare（S1）；编制组五域序列（S3）。统计期：2025.07—2026.06；五域与国家公开口径独立呈现。}
\Box{15}{211}{87}{\fontsize{10.5}{14.2}\selectfont\color{teal}\sffamily \bfseries ① 项目影响向主要协作节点汇聚}
\Box{15}{217}{87}{\fontsize{10.1}{15.6}\selectfont\color{ink}在已统计的 29 国中，美国项目影响力为 144.78 万 OpenRank，中国为 53.07 万，构成两个主要节点。项目归属反映影响汇聚的位置，贡献则可能来自多个国家；理解这种集中格局，需要同时看跨境协作。}
\Box{108}{211}{87}{\fontsize{10.5}{14.2}\selectfont\color{teal}\sffamily \bfseries ② 合作网络需要连接多样化节点}
\Box{108}{217}{87}{\fontsize{10.1}{15.6}\selectfont\color{ink}荷兰、德国和英国分别为 5.51 万、3.16 万和 2.32 万，南非也进入前二十。主要节点之外，各地仍有值得关注的项目和维护团队。合作伙伴的选择应结合技术领域与协作需求，不能仅按国家规模排序。}
\Box{15}{245}{87}{\fontsize{10.5}{14.2}\selectfont\color{teal}\sffamily \bfseries ③ 从广泛参与走向持续组织能力}
\Box{15}{251}{87}{\fontsize{10.1}{15.6}\selectfont\color{ink}同一组 29 国中，美中合计占活跃人数 34.1\%、开发者贡献 37.1\%，项目影响占 89.1\%，三项份额分别计算。参与分布与项目影响分布存在明显差异。这提示公共投入还需支持项目组织、共同维护与协作接续。}
\Box{108}{245}{87}{\fontsize{10.5}{14.2}\selectfont\color{teal}\sffamily \bfseries ④ 在全球分工中理解贡献与影响}
\Box{108}{251}{87}{\fontsize{10.1}{15.6}\selectfont\color{ink}五域序列中，欧洲开发者贡献居首、项目影响居第三；美国分别为第三和第一。贡献来源与项目归属呈现不同次序。结合 2.2 的跨境连接，评价生态应同时关注本地供给与参与全球共建的能力。}
\end{canvas}
\begin{canvas}\Header{第二章 · 全球开源生态地图}{teal}\Footer{02}
\Box{15}{27}{180}{\fontsize{21}{27}\selectfont\color{ink}\sffamily \bfseries 2.3 项目影响的增长与长期情景}
\Box{15}{40}{180}{\fontsize{9.7}{13.5}\selectfont\color{muted}\sffamily 历史曲线呈现规模与增长节奏，2030 情景展示不同假设下的可能路径。五域按报告年度、全球与中美按自然年分别观察；相对位置与绝对差距共同解释格局变化。}
\Box{15}{53}{180}{\includegraphics[width=180mm]{.build/assets/02-project-influence-evolution.pdf}}
\Box{15}{199}{180}{\fontsize{7.8}{10.8}\selectfont\color{muted}\sffamily 图 2-9  项目影响历史与 2030 条件情景。来源：OpenShare 自然年输出（S2）、编制组五域序列（S3）；情景由本报告测算。实线为历史，虚线为情景；区间不是概率或置信区间。}
\Box{15}{211}{87}{\fontsize{10.5}{14.2}\selectfont\color{teal}\sffamily \bfseries ① 快速增长仍需持续建设支撑}
\Box{15}{217}{87}{\fontsize{10.1}{15.6}\selectfont\color{ink}五域历史中，中国项目影响增长 159.9\%。按增速延续或逐年放缓推演，2030 年中国为 130—187 万，美国为 205—229 万。两种情景均指向中国相对位置上升；路径能否延续，仍需跟踪项目维护与协作的稳定性。}
\Box{108}{211}{87}{\fontsize{10.5}{14.2}\selectfont\color{teal}\sffamily \bfseries ② 中国项目影响在全球扩展中上升}
\Box{108}{217}{87}{\fontsize{10.1}{15.6}\selectfont\color{ink}2021—2025 自然年，全球、中国、美国读数分别增长 55.0\%、135.0\%、37.2\%。中国增幅高于全球，显示其在已观测项目网络中的影响扩展较快。该序列与五域覆盖不同，也不足以单独识别 AI 对增长的作用。}
\Box{15}{245}{87}{\fontsize{10.5}{14.2}\selectfont\color{teal}\sffamily \bfseries ③ 相对追近与绝对差距需要并看}
\Box{15}{251}{87}{\fontsize{10.1}{15.6}\selectfont\color{ink}五域序列中，中国相当于美国的比例由 20.7\% 升至 41.2\%；同期绝对差距由 106.5 万降至 103.0 万 OpenRank。较快增速与较大存量差距并存，提示应同时支持成长项目和长期积累，持续观察格局变化。}
\Box{108}{245}{87}{\fontsize{10.5}{14.2}\selectfont\color{teal}\sffamily \bfseries ④ 为全球南方拓展持续共建路径}
\Box{108}{251}{87}{\fontsize{10.1}{15.6}\selectfont\color{ink}南方六国项目影响增长 42.3\%，印度增长 47.9\%，均高于美国，但基数较小。AI 普惠合作可支持本地化、共同维护与区域公共项目，并在未来 6—12 个月跟踪贡献是否被采纳、维护是否持续，使参与机会转为公共能力。}
\end{canvas}
\begin{canvas}\Header{第二章 · 全球开源生态地图}{teal}\Footer{03}
\Box{15}{27}{180}{\fontsize{21}{27}\selectfont\color{ink}\sffamily \bfseries 2.3 中国项目：结构与增长}
\Box{15}{40}{180}{\fontsize{9.7}{13.5}\selectfont\color{muted}\sffamily 从项目结构、历史增长和参与位置观察中国的建设基础。30 项项目标签用于比较内部结构，五域年度序列用于观察整体演变；两者独立呈现，不拼接为同一总体。}
\Box{15}{53}{180}{\includegraphics[width=180mm]{.build/assets/03-china-project-structure.pdf}}
\Box{15}{199}{180}{\fontsize{7.8}{10.8}\selectfont\color{muted}\sffamily 图 2-10  中国项目的样本结构、五年增长与参与位置。来源：编制组归档（S3）。①②④为 30 项标签样本，窗口待核；③为 2022—2026 报告年度五域序列。两类口径独立。}
\Box{15}{211}{87}{\fontsize{10.5}{14.2}\selectfont\color{teal}\sffamily \bfseries ① 系统与计算基础软件构成重要支点}
\Box{15}{217}{87}{\fontsize{10.1}{15.6}\selectfont\color{ink}30 项归档样本中，OpenHarmony 与 openEuler 位居前两名，CANN、Ascend、MindSpore 随后，系统与计算基础软件较为突出。长期支持应落实到维护人员、版本接续和兼容验证，保障开放成果能够持续使用。}
\Box{108}{211}{87}{\fontsize{10.5}{14.2}\selectfont\color{teal}\sffamily \bfseries ② 保障关键项目，也培育多样化供给}
\Box{108}{217}{87}{\fontsize{10.1}{15.6}\selectfont\color{ink}原表读数中，前两项占 53.7\%，前十项占 84.3\%。基金会可兼顾关键项目的稳定维护和成长项目的孵化，连接不同组织与技术方向。这里反映的是标签样本结构；标签可能重叠，不外推为全国集中度。}
\Box{15}{245}{87}{\fontsize{10.5}{14.2}\selectfont\color{teal}\sffamily \bfseries ③ 五年增长需要转化为持续积累}
\Box{15}{251}{87}{\fontsize{10.1}{15.6}\selectfont\color{ink}五域序列中，中国项目影响由 27.74 万增至 72.10 万 OpenRank，增长 159.9\%，2026 年相当于美国的 41.2\%。2023 年增幅约 53.7\%，随后三年为 17.7\%—21.1\%。扩张持续，但节奏并不均匀，长期支持应关注版本接续、维护者留存和跨组织协作。}
\Box{108}{245}{87}{\fontsize{10.5}{14.2}\selectfont\color{teal}\sffamily \bfseries ④ 参与规模与网络位置提供互补信息}
\Box{108}{251}{87}{\fontsize{10.1}{15.6}\selectfont\color{ink}openGauss 的影响位次为第 9、活跃位次为第 19；PaddlePaddle 分别为第 6 和第 10。人数与网络位置呈现不同次序，提示项目支持还需结合协作结构、采用情况和维护需求；两者相除不能解释为个人效率。}
\end{canvas}
\begin{canvas}\Header{附录 1 · 2.3 指标与方法}{teal}\Footer{A1}
\Box{15}{28}{180}{\fontsize{20}{27}\selectfont\color{ink}\sffamily \bfseries 附录 1｜指标、时期与情景方法}
\Box{15}{47}{180}{\fontsize{11.5}{17}\selectfont\color{teal}\sffamily \bfseries 1  三个视角，观察不同对象}
\Box{15}{57}{180}{\fontsize{9.8}{15.5}\selectfont\color{ink}活跃人数回答“有多少人参与”，开发者贡献回答“贡献来自哪里”，项目影响回答“影响汇聚到哪些项目”。本节项目影响采用 OpenRank，按项目国家标签汇总；国家标签不是法律所有权，也不说明贡献者国籍。指标不能直接替代技术质量、收入或公共服务成效。}
\Box{15}{88}{180}{\fontsize{11.5}{17}\selectfont\color{teal}\sffamily \bfseries 2  同集合比较与五域序列}
\Box{15}{98}{180}{\fontsize{9.7}{15.5}\selectfont\color{ink}图 2-8③固定项目影响输出覆盖的 29 个国家，对活跃人数、开发者贡献和项目影响分别求和，再计算美国、中国、其余 27 国份额。三个分母分别为 2,181.30 万人、1,147.66 万 OpenRank 和 222.08 万 OpenRank；不是全球总体分母。项目标签覆盖不足会影响判断。份额独立舍入，合计可能有尾差。}
\Box{15}{129}{180}{\fontsize{9.7}{15.5}\selectfont\color{ink}五域采用编制组独立综合序列，不与国家公开读数混排或求差归因。南方六国指巴西、印度尼西亚、越南、南非、菲律宾、埃及，印度单列；六国不代表全球南方总体。欧洲成员范围和底层聚合映射仍待完整核验。}
\Box{15}{157}{180}{\fontsize{11.5}{17}\selectfont\color{teal}\sffamily \bfseries 3  历史与 2030 情景}
\Box{15}{167}{180}{\fontsize{9.7}{15.5}\selectfont\color{ink}当期统计为 2025-07-01 至 2026-06-30。五域历史为 2022—2026 报告年度，推演始于 2027；全球及中美趋势另用 2021—2025 自然年，推演始于 2026。自然年图中的 2026 点也是情景值，不是已观测当期读数。}
\Box{15}{194}{180}{\fontsize{9.7}{15.5}\selectfont\color{ink}历史年均增速 g =（末年值 ÷ 首年值）的四次方根 − 1。延续情景每年按 g 复利；放缓情景第 h 年增速为 g × 0.8 的 h 次方，再逐年复利。0.8 为敏感性参数，两条路径不赋予发生概率，未估计 AI 的因果效应。年度覆盖及指标可比是解释成立的条件。}
\Box{15}{225}{180}{\fontsize{11.5}{17}\selectfont\color{teal}\sffamily \bfseries 4  中国项目样本如何阅读}
\Box{15}{235}{180}{\fontsize{9.7}{15.5}\selectfont\color{ink}原表含 30 项目标签，影响读数之和为 176,313.03。集中度按原表直接计算，未主张标签间已去重；活跃位次按同表人数重排，项目间人数也不相加。原表统计窗口未单列，故仅作为归档结构样本。图 2-10③另用五域中国序列，不能与 30 项样本求和或计算覆盖率；2022—2026 读数依次为 27.74、42.64、50.20、60.80、72.10 万 OpenRank。}
\end{canvas}
\begin{canvas}\Header{附录 1 · 来源与评审记录}{teal}\Footer{A2}
\Box{15}{28}{180}{\fontsize{20}{27}\selectfont\color{ink}\sffamily \bfseries 附录 1｜数据出处与评审记录}
\Box{15}{45}{180}{\fontsize{11.5}{17}\selectfont\color{teal}\sffamily \bfseries 2030 情景读数 · 单位：万 OpenRank}
\Box{15}{55}{180}{\fontsize{9.4}{14.8}\selectfont\color{ink}五域（报告年度）：美国 205.36—228.58；中国 129.92—187.42；欧洲 22.05—24.94；南方六国 0.74—0.85；印度 0.52—0.61。自然年序列：全球 370.25—473.84；美国 233.51—279.31；中国 100.32—160.28。均按“放缓—延续”排列。}
\Box{15}{83}{180}{\fontsize{11.5}{17}\selectfont\color{teal}\sffamily \bfseries 可点击的数据与项目出处}
\Box{15}{94.0}{180}{\sffamily\fontsize{8.6}{12}\selectfont\textcolor{muted}{S1}\hspace{2mm}\href{https://selfoss.open-digger.cn/openshare/overview_half_year/meta.json}{报告年度：影响、贡献、开发者及统计期}}
\Box{15}{99.6}{180}{\sffamily\fontsize{8.6}{12}\selectfont\textcolor{muted}{S1a}\hspace{2mm}\href{https://selfoss.open-digger.cn/openshare/overview_half_year/influence.json}{当期项目影响原始 JSON}}
\Box{15}{105.2}{180}{\sffamily\fontsize{8.6}{12}\selectfont\textcolor{muted}{S1b}\hspace{2mm}\href{https://selfoss.open-digger.cn/openshare/overview_half_year/contribution.json}{当期开发者贡献原始 JSON}}
\Box{15}{110.8}{180}{\sffamily\fontsize{8.6}{12}\selectfont\textcolor{muted}{S1c}\hspace{2mm}\href{https://selfoss.open-digger.cn/openshare/overview_half_year/developers.json}{当期开发者原始 JSON}}
\Box{15}{116.4}{180}{\sffamily\fontsize{8.6}{12}\selectfont\textcolor{muted}{S2}\hspace{2mm}\href{https://selfoss.open-digger.cn/openshare/overview/influence.json}{自然年项目影响序列}}
\Box{15}{122.0}{180}{\sffamily\fontsize{8.6}{12}\selectfont\textcolor{muted}{S3}\hspace{2mm}\href{https://github.com/will-ww/2026-GOSD-Report/tree/497478a19e92c676c225893a395738e747a71c8b/reports/round2-v2.0/source/chapters/02/modules/02/data/23}{五域序列、中国 30 项目原表及派生表}}
\Box{15}{127.6}{180}{\sffamily\fontsize{8.6}{12}\selectfont\textcolor{muted}{S4}\hspace{2mm}\href{https://github.com/X-lab2017/open-digger/blob/5c2ffc4f5b3dd37ebcb20514c7abfee6853bd830/src/dataProduction/openshareOverview.ts}{生产口径与公开数值调整说明（固定代码）}}
\Box{15}{133.2}{180}{\sffamily\fontsize{8.6}{12}\selectfont\textcolor{muted}{S5}\hspace{2mm}\href{https://github.com/X-lab2017/open-digger/blob/master/skills/open-digger-data-retrieval/SKILL.md}{OpenDigger 指标与获取方法}}
\Box{15}{172}{180}{\fontsize{11.5}{17}\selectfont\color{teal}\sffamily \bfseries 本版修改与主张衔接}
\Box{15}{182}{180}{\fontsize{9.5}{15}\selectfont\color{ink}v0.3（2026-11-10）：图 2-10 删除技术用途分层，改为五域中国序列的五年增长；新增分析置于③，参与与影响位次移至④，文字同步调整。保留地图标注修正。数据、统计窗口与情景方法沿用前稿。摘要与第六章可吸收的方向：广泛参与需要持续组织能力；中国贡献应放在全球接续中理解；普惠投入应支持可持续公共能力。}
\Box{15}{214}{180}{\fontsize{11.5}{17}\selectfont\color{teal}\sffamily \bfseries 待核查项与评审问题}
\Box{15}{224}{180}{\fontsize{9.4}{14.8}\selectfont\color{ink}待核查：① 公开生产代码含估算与年度调整，其固定版本与所取 JSON 的对应关系；② 五域的成员集合与聚合映射；③ 中国 30 项目统计窗口、活跃人数定义与标签去重边界。这些问题限制解释范围，不以新的猜测补值。原始 JSON、CSV、派生计算和检索日期随源稿保留。}
\Box{15}{254}{180}{\fontsize{9.3}{14.6}\selectfont\color{ink}请重点评审：地图标注是否清晰；12 个分析点是否兼顾证据、解释与决策意义，能否衔接后续人才、资产和协作观察。依据：报告宪法批准基线 fff392e；当前材料提交 497478a。}
\end{canvas}\end{document}
